\documentclass[11pt,a4paper]{article}
\usepackage[T1]{fontenc}
\usepackage[utf8]{inputenc}
\usepackage[margin=2.4cm]{geometry}
\usepackage{booktabs}
\usepackage{array}
\newcommand{\code}[1]{\texttt{#1}}

\title{Étude de la nouvelle codation des couleurs\\[2pt]
\large Simulation des méthodes sur les données réelles Rubicon}
\author{Suite Rubicon}
\date{7 septembre 2026}

\begin{document}
\maketitle

\section{Contexte et objet de l'étude}

Chaque produit du catalogue porte un code de la forme
\code{\{MODÈLE\}-\{COULEURS\}/\{MÉTAL\}}, par exemple
\code{B012-AM+PER/W}~: le modèle \code{B012}, monté en or blanc
(\code{W}), serti d'améthyste et de péridot. Le segment COULEURS juxtapose
les codes des pierres, joints par \code{+}, la pierre centrale d'abord puis
les autres par poids décroissant.

Ce segment n'obéit aujourd'hui à aucune règle stricte~: la nuance de la
pierre (sa \emph{shade}~—- \code{AA}, \code{PL Pink Light}\ldots) y est
tantôt écrite, tantôt omise, et la forme de taille (\emph{shape}~—- ronde,
poire\ldots) n'y figure jamais. Il en résulte des codes hétérogènes et un
risque de collision~: deux produits aux pierres différentes peuvent porter
le même code. L'objectif d'une nouvelle codation est double et
contradictoire~—- des codes \emph{les plus courts et lisibles possible},
mais \emph{sans perte d'information ni collision}. Plusieurs grammaires
candidates ont été proposées~; cette étude les simule sur l'usage réel pour
les départager. Elle traite aussi deux questions connexes~: le codage des
boucles d'oreilles vendues à l'unité, et l'existence éventuelle d'un
standard de place.

Contrainte de continuité posée d'emblée~: quelle que soit la codation
retenue, les codes historiques ne seront jamais réécrits~—- ils restent
consultables et cherchables, avec une table de correspondance entre ancien
et nouveau système (voir §\ref{sec:continuite}).

\section{Données de l'étude}

L'étude porte sur l'usage réel des pierres dans les compositions de produits~:

\begin{center}
\begin{tabular}{lr}
\toprule
Lignes de composition (occurrences pondérantes) & 247\,898 \\
Triplets distincts (type, shade, forme) utilisés & 2\,062 \\
Types de pierre au dictionnaire & 112 \\
Shades au dictionnaire & 115 \\
Formes au dictionnaire & 178 \\
\bottomrule
\end{tabular}
\end{center}

\noindent Les occurrences sont extraites des \emph{compositions r\'eelles}~:
chaque ligne de pierre pos\'ee sur un produit compte, pond\'er\'ee par sa
fr\'equence. Les \emph{shades} disposent bien d'un code au dictionnaire
(115 codes, longueur moyenne 2,5, maximum 5~: \code{A}, \code{AA},
\code{AAA}, \code{W1}, \code{PL}, \code{LC}, \code{T2}\ldots)~; la
dysharmonie actuelle vient donc de leur \emph{usage} dans les codes
produits (tant\^ot \'ecrits, tant\^ot omis, sans r\`egle), pas d'une absence
de codes. Les m\'ethodes B et C sont simul\'ees avec ces codes r\'eels.

Les statistiques sont pondérées par le nombre d'occurrences~: un code porté
par 30\,000 lignes pèse 30\,000 fois plus qu'un code rare. Les valeurs par
défaut (shade et forme) sont calculées \emph{par type}~: la shade par défaut
du diamant est \code{White1~=~Si-P1}, sa forme par défaut \code{RDWFC},
celle des pierres de couleur \code{RD}.

\section{Grammaires candidates simulées}

\begin{description}
\item[A — type seul] Le jeton est le code du type~; shade et forme
  disparaissent du code.
\item[B — \code{TYPE.SHADE}] La shade n'est écrite que si elle diffère de la
  shade par défaut du type.
\item[C — \code{TYPE.SHADE.SHAPE}] Comme B, plus la forme si elle diffère de
  la forme par défaut (\code{..} quand la shade est par défaut mais pas la
  forme).
\item[D — \code{PPSSHH}] Recodage de tout le dictionnaire en 2 caractères
  (types, shades, formes), valeurs par défaut omises. Trois variantes pour
  le cas « forme sans shade » (\code{PPHH})~: alphabets disjoints
  (aucun séparateur), point de remplacement (\code{PP.HH}), ou shade
  obligatoire dès que la forme est écrite (\code{PPSSHH}).
\item[E — suffixe \code{[\,]}] Non exclusif~: un discriminant optionnel
  ajouté seulement en cas de collision résiduelle, quelle que soit la
  méthode porteuse.
\end{description}

\section{Résultats}

\begin{center}
\begin{tabular}{l r r r r l}
\toprule
Méthode & Moy. & Max & Variance & Écart-type & Information perdue \\
\midrule
A — type seul            & 2,85 &  5 & 0,58 & 0,76 & shade~: 97,6\,\% des occurrences ambiguës \\
B — \code{TYPE.SHADE}    & 4,03 &  9 & 2,17 & 1,47 & forme (comme aujourd'hui) \\
C — \code{TYPE.SHADE.SHAPE} & 5,04 & 15 & 5,44 & 2,33 & aucune \\
D — disjoint             & \textbf{3,21} & \textbf{6} & 1,53 & 1,24 & aucune \\
D — point                & 3,36 &  6 & 1,89 & 1,38 & aucune \\
D — shade obligatoire    & 3,50 &  6 & 2,50 & 1,58 & aucune \\
\bottomrule
\end{tabular}
\end{center}

Chiffres utiles à la lecture~:
38,7\,\% des occurrences portent une shade non-défaut (donc écrite)~;
21,9\,\% une forme non-défaut~;
le cas \code{PPHH} (forme écrite, shade par défaut) représente
\textbf{14,6\,\%} des occurrences — c'est lui que les trois variantes de D
départagent. La distribution actuelle des longueurs de codes de type~:
32\,\% en 2 caractères, 55\,\% en 3, 12\,\% en 4--5.

\subsection{Deux niveaux de lecture~: jeton pierre et code produit}

Les ambiguïtés du tableau ci-dessus se lisent au niveau du \emph{jeton
pierre}~: 97,6\,\% signifie qu'en lisant un jeton de la méthode A on ne
peut pas, dans 97,6\,\% des occurrences, retrouver la shade de la pierre.
C'est une mesure de perte d'information, pas un taux de doublons.

La question pratique est ailleurs~: deux \emph{produits} du même modèle et
du même métal peuvent-ils porter le même code~? Mesuré sur les
45\,895 produits disposant d'une composition de pierres, en groupant par
(modèle, métal, segment couleurs)~:

\begin{center}
\begin{tabular}{l r r}
\toprule
Grammaire & Produits en collision & Groupes \\
\midrule
A — type seul (vs shade)         & 13\,907 (30,3\,\%) & 4\,989 \\
A — type seul (vs shade+forme)   & 14\,780 (32,2\,\%) & 5\,373 \\
B — type+shade (vs forme)        &  1\,846 (4,0\,\%)  &   799 \\
C / D — triplet complet          & 0 & — \\
\bottomrule
\end{tabular}
\end{center}

La médiane des modèles n'a qu'\emph{un seul} code couleurs (moyenne 3,0,
maximum 79)~: l'intuition «~un modèle donné collisionne rarement~» est
vraie pour le modèle médian, mais les modèles multi-variantes concentrent
la masse~—- \textbf{30\,\% des produits} partageraient leur code avec un
frère de shades différentes sous la méthode A. La méthode B ramène ce
taux à 4\,\% (variantes de forme seule), C et D à zéro.

Reste un dernier chiffre, indépendant de la grammaire~: \textbf{12,3\,\%}
des produits (5\,661, en 2\,489 groupes) ont un frère aux pierres
\emph{identiques} à shade et forme près~—- ils ne diffèrent que par le
nombre, la taille ou la disposition des pierres, qu'aucun code couleurs
n'encode. C'est précisément le domaine du suffixe \code{[\,]}
(§\ref{sec:collisions})~: il n'est pas anecdotique et devra être
spécifié.

\subsection{Exemples concrets}

Shades codées en exemple~: \code{1C} = Light Ceylon. \code{..} marque la
shade par défaut omise en C.

\begin{center}
\small
\begin{tabular}{l l l l l l}
\toprule
Pierre réelle & A & B & C & D-disjoint & D-point \\
\midrule
Diamant blanc défaut (RDWFC)      & \code{DTS}   & \code{DTS}    & \code{DTS}      & \code{DI}   & \code{DI} \\
Saphir Light Ceylon, RD           & \code{SA}    & \code{SA.LC}  & \code{SA.LC}    & \code{SA1C} & \code{SA1C} \\
Améthyste défaut, poire (PS)      & \code{AME}   & \code{AME}    & \code{AME..PS}  & \code{AMPS} & \code{AM.PS} \\
Tourmaline rose défaut, RD        & \code{PTOUR} & \code{PTOUR}  & \code{PTOUR}    & \code{PT}   & \code{PT} \\
\bottomrule
\end{tabular}
\end{center}

En A, le saphir Light Ceylon et le saphir rose s'écrivent tous deux
\code{SA}~: c'est la collision type~—~le saphir compte 32 shades utilisées,
le diamant 23.

\section{Fréquences et attribution des codes courts}

Top 20 des types par usage, avec une \emph{proposition} d'allocation
2~caractères (à valider~—- aucun doublon dans la proposition)~:

\begin{center}
\small
\begin{tabular}{l l r r r l}
\toprule
Code actuel & Type & Occurrences & Part & Shades utilisées & Proposition \\
\midrule
DTS   & Diamond          & 56\,261 & 22,7\,\% & 23 & \code{DI} \\
SA    & Sapphire         & 48\,746 & 19,7\,\% & 32 & \code{SA} \\
BT    & Blue Topaz       & 16\,894 &  6,8\,\% &  4 & \code{BT} \\
TSAV  & Tsavorite        & 14\,938 &  6,0\,\% &  4 & \code{TS} \\
AME   & Amethyst         & 13\,158 &  5,3\,\% &  6 & \code{AM} \\
CIT   & Citrine          & 10\,371 &  4,2\,\% &  4 & \code{CI} \\
IOL   & Iolite           & 10\,051 &  4,1\,\% &  3 & \code{IO} \\
PTOUR & Pink Tourmaline  &  9\,656 &  3,9\,\% &  2 & \code{PT} \\
PER   & Peridot          &  9\,402 &  3,8\,\% &  3 & \code{PE} \\
RHO   & Rhodolite        &  9\,224 &  3,7\,\% &  1 & \code{RH} \\
MON   & Moon Stone       &  4\,723 &  1,9\,\% & 10 & \code{MO} \\
RDF   & Rose de France   &  4\,428 &  1,8\,\% &  0 & \code{RF} \\
GA    & Garnet           &  3\,827 &  1,5\,\% &  3 & \code{GA} \\
PRL   & Pearl            &  3\,653 &  1,5\,\% & 12 & \code{PL} \\
GQ    & Green Quartz     &  3\,136 &  1,3\,\% &  1 & \code{GQ} \\
GTOUR & Green Tourmaline &  2\,865 &  1,2\,\% &  4 & \code{GT} \\
PQA   & Pink Quartz      &  2\,375 &  1,0\,\% &  6 & \code{PQ} \\
SMQ   & Smokey Quartz    &  2\,325 &  0,9\,\% &  3 & \code{SQ} \\
RU    & Ruby             &  2\,104 &  0,8\,\% &  6 & \code{RU} \\
LAQ   & Lemon Quartz     &  1\,774 &  0,7\,\% &  1 & \code{LQ} \\
\bottomrule
\end{tabular}
\end{center}

Ces 20 types couvrent 92\,\% des occurrences~; les 92 types restants se
partagent les 8\,\% et peuvent recevoir des codes moins mnémoniques sans
douleur. Un outil d'allocation automatique par fréquence (avec veto manuel)
est facile à fournir.

\section{Un préalable~: les types porteurs de couleur}

Le dictionnaire des types mélange trois familles~: les \emph{espèces pures}
(\code{SA} Sapphire, \code{TOU} Tourmaline, \code{TOP} Topaz), les
\emph{noms de variété} consacrés par la gemmologie (\code{AME} Amethyst,
\code{CIT} Citrine, \code{RHO} Rhodolite, \code{TSAV} Tsavorite~—- pas un
problème~: personne ne dit «~quartz violet~»), et \textbf{14 composés
ad hoc couleur+espèce}~: \code{BT} Blue Topaz, \code{PTOUR} Pink
Tourmaline, \code{GTOUR}, \code{GQ} Green Quartz, \code{PQA}, \code{SMQ},
\code{LAQ}, \code{WT}, \code{WQA}, \code{CHP}, \code{RUT}, \code{OLQ},
\code{BKTOU}, \code{BTOUR}. Ces composés portent \textbf{40\,184
occurrences (16\,\% de l'usage)}, alors que les espèces de base
coexistent, quasi inutilisées (\code{TOU}~: 14 occurrences, \code{TOP}~:
19). Le même concept est donc encodé sur deux axes différents~: le saphir
rose s'écrit \code{SA}+shade \code{PL}, la tourmaline rose s'écrit
\code{PTOUR} tout court.

Le cumul redouté «~type-couleur + shade~» existe~: \textbf{19\,402
lignes}. Mais 98,6\,\% d'entre elles portent une shade de \emph{grade}
(\code{BT}+\code{AA} = topaze bleue, ton moyen~—- 16\,891 lignes pour la
seule topaze bleue, 2\,169 pour le quartz fumé)~: sémantiquement cohérent,
structurellement impur. Les vraies redondances couleur-sur-couleur sont
marginales et à nettoyer~: 3 lignes \code{PQA}+\code{PK} («~quartz rose~»
rose). L'axe shade mélange d'ailleurs lui-même trois notions~: la
\emph{teinte} (\code{PK}, \code{CD} chrome), le \emph{grade}
(\code{A}/\code{AA}/\code{AAA}/\code{2}/\code{3}/\code{C}) et le
\emph{phénomène} (\code{CDY} Cloudy, \code{ML} Milky, \code{SR} Star).

La question de fond est donc~: \emph{sur quel axe la couleur doit-elle
vivre~?} Deux architectures s'opposent~:

\begin{enumerate}
\item \textbf{Tout en article}~: chaque couple espèce+couleur devient un
  type à part entière (généraliser ce que \code{BT} et \code{PTOUR} font
  déjà), avec au besoin une convention de lettres pour les couleurs dans
  les codes de types.
\item \textbf{Tout en shade éclatée}~: la couleur reste dans l'axe shade,
  mais celui-ci se scinde en trois sous-axes~—- teinte, grade,
  phénomène~—- chacun avec un code d'un caractère.
\end{enumerate}

\subsection{Arbitrage recommandé~: l'article porte la couleur, la qualité
s'éclate}
\label{sec:axes}

Les données tranchent~: les deux options règlent chacune une moitié du
problème et se combinent.

Côté shade, les 115 codes se répartissent en \textbf{7 grades, 9
phénomènes et 99 «~couleurs et autres~»}~—- et ces 99 sont pour beaucoup
des fusions teinte$\times$ton (\code{PL}~= Pink \emph{Light}, \code{PM}~=
Pink \emph{Medium}, \code{LC}/\code{OC}~= ceylon \emph{light/open})~: la
liste des teintes pures est courte ($\sim$15--25), une lettre y suffit.
Côté type, généraliser l'article espèce+couleur porte le dictionnaire de
112 à \textbf{277 entrées} en usage réel (gros contributeurs~: diamants
\code{DTS+W1}/\code{T2}, saphirs \code{PL}/\code{LC}/\code{OC}/\code{PM})
~—- confortable dans un espace 2 caractères (676 combinaisons).

L'hybride retenu~:
\begin{itemize}
\item \textbf{l'identité dans le type}~: la couleur monte dans le type,
  qui devient l'\emph{article commercial} (Pink Sapphire, Blue Topaz,
  White Diamond SI)~—- ce que le métier nomme et price, et le modèle
  même d'Emasur~; l'asymétrie \code{SA}+\code{PL} / \code{PTOUR}
  disparaît~;
\item \textbf{la qualité dans la shade, scindée}~: grade sur un
  caractère (chiffre~—- 7 valeurs) et phénomène sur une lettre (9
  valeurs)~; plus aucune teinte dans la shade \emph{par construction},
  la redondance couleur-sur-couleur devient impossible~;
\item la grammaire D devient \code{PP}~(article) + \code{G}~(grade,
  chiffre) + \code{P}~(phénomène, lettre) + \code{HH}~(forme)~—-
  alphabets disjoints par construction, codes encore plus courts
  puisque la teinte a disparu dans \code{PP}~;
\item la convention de lettres-couleurs dans les codes d'articles reste
  une \emph{préférence} de l'outil d'allocation (systématisable en
  3 caractères, trop à l'étroit en 2 pour 277 articles).
\end{itemize}

Coût assumé~: le dictionnaire d'articles est vivant (une nouvelle couleur
de quartz = une nouvelle ligne)~—- d'où l'écran de curation et l'outil
d'allocation par fréquence. En échange, la saisie garde deux axes
principaux (article, qualité) au lieu de quatre menus. En attendant le
nouveau système, un garde-fou reste possible sur l'ancien~: interdire une
shade de teinte sur un type déjà porteur de couleur (les 3 lignes
\code{PQA}+\code{PK} sont à nettoyer).

\subsection{Exemple fil rouge~: une même bague sous chaque grammaire}

Bague \code{B012} en or blanc~—- centre \textbf{saphir ceylon clair
taillé poire}, entourage \textbf{topaze bleue grade AA ronde}, pavage
\textbf{diamants blancs baguette}. Codes d'exemple pour les schémas
recodés~: \code{LC}$\to$\code{1C}, \code{W1}$\to$\code{1W}, grade
\code{AA}$\to$\code{2}, articles \code{SC}~Ceylon Sapphire,
\code{BT}~Blue Topaz, \code{DI}~White Diamond. Défauts~: shade
\code{PL} et forme \code{RD} pour le saphir, shade \code{A} et forme
\code{RD} pour la topaze, shade \code{W1} et forme \code{RDWFC} pour le
diamant.

\begin{center}
\small
\begin{tabular}{l l r l}
\toprule
Grammaire & Segment couleurs & Long. & Information perdue \\
\midrule
Aujourd'hui (sans règle) & \code{SA+BTAA+D} \emph{ou} \code{SALC+BT+DTS}\ldots & — & variable, non prédictible \\
A — type seul            & \code{SA+BT+DTS}            &  9 & teintes, grades, formes \\
B — \code{TYPE.SHADE}    & \code{SA.LC+BT.AA+DTS}      & 15 & formes \\
C — complet              & \code{SA.LC.PS+BT.AA+DTS..BG} & 22 & aucune \\
D — disjoint             & \code{SA1CPS+BT2A+DIBG}     & 16 & aucune \\
D — point                & \code{SA1CPS+BT2A+DI.BG}    & 17 & aucune \\
D — shade obligatoire    & \code{SA1CPS+BT2A+DI1WBG}   & 18 & aucune \\
\textbf{Hybride (retenu)} & \code{SC1PS+BT2+DIBG}     & \textbf{14} & aucune \\
\bottomrule
\end{tabular}
\end{center}

Lecture de l'hybride~: \code{SC1PS} = article saphir ceylon, grade
1~(light), forme poire~; \code{BT2} = topaze bleue, grade 2~(AA), forme
par défaut~; \code{DIBG} = diamant blanc, grade par défaut, baguette. La
teinte ayant rejoint l'article, c'est le segment le plus court qui ne
perde \emph{aucune} information~—- et chaque bloc se découpe sans
séparateur (article~= lettres, grade~= chiffre, forme~= lettres).

\section{La question \code{PPHH} et la variante recommandée}

Le cas « forme écrite, shade par défaut » touche 14,6\,\% des
occurrences. Les trois options~:

\begin{enumerate}
\item \textbf{Alphabets disjoints (recommandée)}~: les codes de shade
  commencent par un \emph{chiffre} (ex.~\code{1C}, \code{2A}), les codes de
  forme par une \emph{lettre}. $10 \times 36 = 360$ combinaisons pour 115
  shades, largement suffisant. Les codes de shade actuels \'etant
  surtout alphab\'etiques (\code{A}, \code{PL}, \code{GYAAA}\ldots, avec
  d\'ej\`a \code{2} et \code{3} num\'eriques), D implique de recoder les 115
  shades~—- table de correspondance triviale. Le code \code{SA1CPS} se découpe alors sans
  ambiguïté ni séparateur~: moyenne 3,21, maximum 6.
\item \textbf{Point de remplacement} (\code{AM.PS})~: lisible, +0,15 de
  moyenne, réintroduit un séparateur dans 14,6\,\% des codes.
\item \textbf{Shade obligatoire avec la forme}~: grammaire la plus simple
  mais la plus longue (+0,29) et oblige à écrire une information par
  défaut.
\end{enumerate}

\section{Collisions résiduelles et suffixe \code{[\,]}}
\label{sec:collisions}

Avec D (ou C), le triplet complet est encodé~: il n'y a plus de collision
\emph{de pierre}. Le suffixe \code{[\,]} reste utile à un autre niveau~: deux
\emph{produits} du même modèle portant les mêmes pierres en composition
différente (nombre, disposition). Mesuré sur les données~: 12,3\,\% des produits sont dans ce cas. Il doit rester optionnel~;
l'unicité du code produit est déjà contrôlée par le logiciel (contrainte
d'unicité + suggestion de code calculée).

\section{Recommandation}

\textbf{Méthode D, variante alphabets disjoints.} C'est la seule qui soit à
la fois courte (3,21 en moyenne, jamais plus de 6), complète (shade et forme
encodées quand elles s'écartent du défaut), bornée (largeur fixe par bloc,
lisible par machine sans séparateur) et harmonieuse (une seule règle pour
tout le dictionnaire). Combinée à l'arbitrage des axes (§\ref{sec:axes})~—- l'article porte
la couleur, la qualité s'éclate en grade et phénomène~—- elle donne la
grammaire cible \code{PP\,G\,P\,HH}. Son coût~: recoder les dictionnaires
et assumer une table de correspondance — coût déjà
provisionné, voir §\ref{sec:continuite}. Si ce recodage devait être refusé,
la méthode B (\code{TYPE.SHADE} avec défauts) est le meilleur compromis sans
recodage~: elle règle l'incohérence d'usage des shades (le problème réellement
vécu) pour +1,2 caractère en moyenne (4,03, maximum 9), en restant sans
la forme, comme aujourd'hui.

\section{Existe-t-il d\'ej\`a un standard de codes~?}

Pour les \textbf{diamants}, oui~: les 4C du GIA (couleur \code{D--Z},
puret\'e \code{FL--I3}) sont le standard de fait mondial — la shade
\code{W1 = Si-P1} du dictionnaire s'y r\'ef\`ere d\'ej\`a. Pour les
\textbf{pierres de couleur}, non~: il n'existe aucun standard international
de codes. Le GIA propose une \emph{description} teinte/ton/saturation
(ex.~\code{vB 6/5}) — trop verbeuse pour servir de code d'identit\'e — et
l'\'echelle commerciale \code{AAA/AA/A}, que nos codes refl\`etent d\'ej\`a,
varie d'un n\'egociant \`a l'autre sans organisme normalisateur (le CIBJO ne
normalise que la nomenclature et la divulgation des traitements). Conclusion~:
rien \`a adopter tel quel~; notre dictionnaire reste propri\'etaire, en
gardant l'ancrage GIA pour le vocabulaire des shades de diamants.

\section{Boucles d'oreilles~: \code{EL} / \code{ER}}

Les boucles d'oreilles sont aujourd'hui référencées uniquement par paire
(catégorie \code{E}). Le besoin nouveau~: distinguer une boucle vendue à
l'unité, oreille gauche ou droite. Deux notations sont en présence~: celle
proposée par Emasur, \code{EXXX*OG} / \code{EXXX*OD} (\emph{oreille
gauche/droite}), et l'alternative \code{EL} / \code{ER} (\emph{earring
left/right}), \code{E} restant la paire. La première mélange français et
anglais dans un dictionnaire jusqu'ici anglophone et introduit un
séparateur inédit (\code{*})~; la seconde s'insère dans la grammaire
existante sans rien ajouter. Recommandation~: \code{EL}/\code{ER}.
Vérification faite en base~: les codes de catégorie \code{EL} et \code{ER}
sont \emph{libres} (l'existant a \code{E} Earring, \code{EB} Earrings
Bottom, \code{BO} Earrings Top) et le mécanisme qui déduit la catégorie du
préfixe alphabétique du code produit les reconnaîtra sans modification.
11\,282 produits commencent par \code{E} aujourd'hui~; rien à migrer, les
codes \code{EL}/\code{ER} ne concernent que les nouvelles références vendues
à l'unité.

\section{Continuité avec l'ancien système}
\label{sec:continuite}

L'exigence~: ne rien écraser, garder l'ancien code consultable, disposer
d'une table de correspondance entre les deux systèmes. \textbf{L'infrastructure existe déjà}~:
chaque produit, modèle, pierre et métal porte un champ \emph{code alternatif}
(avec sa provenance officielle/calculée), l'historique des commandes est
indexé sous les deux notations, la recherche trouve toujours les deux, et
l'affichage se bascule par utilisateur. Les tables de correspondance
(pierre~$\leftrightarrow$~code externe, alias de tokens hérités) sont en
place avec leurs écrans de curation. La nouvelle codation couleur viendra
s'y loger comme un système de notation supplémentaire~: l'ancien code n'est
jamais réécrit (quand un code structuré est appliqué à un produit, l'ancien
est conservé et reste cherchable).

\section{Prochaines étapes}

\begin{enumerate}
\item Valider l'architecture des axes (article porteur de couleur,
  qualité scindée grade/phénomène, §\ref{sec:axes}) et la grammaire
  \code{PP\,G\,P\,HH} (D-disjoint).
\item Attribuer les codes 2 caractères~: outil de proposition automatique
  par fréquence, veto manuel, puis gel du dictionnaire.
\item Fixer la règle d'écriture des shades dans le code produit (les
  codes existent déjà au dictionnaire)~; pour D, les recoder en
  2~caractères à initiale chiffrée.
\item Module de transcription/vérification~: recherche code~$\to$~pierre et
  pierre~$\to$~code (idem shade, forme)~; le contrôle de cohérence d'une
  commande (ordre des pierres~: centre d'abord, puis poids décroissant)
  existe déjà et sera étendu à la nouvelle grammaire.
\item Documentation des calculs~: la règle d'ordre est documentée dans le
  dépôt (\code{rubicon\_addons/pdp\_product/README.md}).
\end{enumerate}

\end{document}
